
from which they are obtained is compatible with the ring structures, using (7.5.1) one finds that
\sgaStableTarget{sga6.e10.eq.7-9-3}@@M0@@
with the pairings @@M1@@ of 4.1 between @@M2@@ and @@M3@@. It follows formally that \emph{if} @@M4@@ \emph{is an element of} @@M5@@ (\emph{resp}. \emph{of} @@M6@@) \emph{such that} @@M7@@ (\emph{resp}. @@M8@@) \emph{is numerically equivalent to zero} (4.5.1), \emph{then the same is true of} @@M9@@. The converse is true when (7.9.2) (or, equivalently, either of the homomorphisms (7.9.1) obtained from it by passage to the associated graded groups) is an epimorphism mod.torsion, - a circumstance that occurs only rather exceptionally.

\sgaStableTarget{sga6.e10.para.7-9-4}\textbf{7.9.4.} It is plausible, at least when @@M10@@ is proper and \emph{smooth}, that the converse in question always holds, i.e., that @@M11@@ (resp. @@M12@@) is numerically equivalent to zero as soon as the same is true of @@M13@@. This would imply that the homomorphisms (7.9.1) induce \emph{monomorphisms} on the corresponding quotient groups by numerical equivalence, @@M14@@ and @@M15@@. At least the preceding discussion implies that the group @@M16@@ is isomorphic to a quotient of a subgroup of @@M17@@ (and in particular has rank at most that of @@M18@@), and likewise for @@M19@@, and also for @@M20@@ itself.

\sgaStableTarget{sga6.e10.sec.7.10}\subsection*{7.10. Compatibility with direct and inverse images}
Let
@@M21@@
be a morphism of noetherian schemes flat over @@M22@@. Denote by @@M23@@ (resp. @@M24@@) the specialization homomorphism for @@M25@@ (resp. @@M26@@), and by @@M27@@ (resp. @@M28@@)
